\section{Examples and limits of coefficient cancellation}\label{sup:cancellation-remarks}
The first two examples illustrate the general algebraic distinction.
For example, over algebraically closed $B$, take a transcendental $x$
and $X=x$, $d=x/2$. Then $V=x/2$ and $B[V]\cap B[X]=B[x]$, although
$\mathcal A_d=B$. If instead $d=2x$, then $P(Z)=Z^2$ gives
$P(V)=P(X)$ without being translation invariant. Thus independence of
$X$ over $B$ leaves mixed relations between $X$ and $\mathbf c$ undetermined,
including the arithmetic nature of $\gamma$.

For connection values, use main-text Theorem~\ref{thm:connection}'s notation.
Linked entries must first be reduced. Contiguity and the determinant give
\[
 \frac{\Gamma(a+1)U(a+1,1,t)}{f_{a+1}(t)}
 -\frac{\Gamma(a)U(a,1,t)}{f_a(t)}
 =-\frac{e^t}{a f_a(t)f_{a+1}(t)}\in B,
\]
and the half-Kummer entry at $(2u,1/2)$ equals $K_0(u)/I_0(u)$.
Higher Ein levels are excluded: their natural tail already satisfies
$\int_x^\infty e^{-v}\log(v/x)\,dv/v
 =\tfrac12(\log x+\gamma)^2+\pi^2/12-\Ein_2(x)$.
No reduction of this connection to $B+B\gamma$ is assumed.

Fischler--Rivoal's Gamma-polynomial connection mechanism
\cite[Theorem~2(iii) and the subsequent discussion on p.~35]{FischlerRivoal2016}
gives asymptotic coefficients at infinity.
They have the form $\Gamma(\alpha^*)P(\gamma)$ over their ring of $G$-values.
Here $\alpha^*=1$ for nonpositive integral $\alpha$.
Gauss's digamma formula and recurrence are classical.
Bundschuh~\cite[Corollary~1]{Bundschuh1979} proved transcendence of
$\psi(a)+\gamma$ for rational nonintegral $a$. See also
Murty--Saradha~\cite[Theorems~1, 3--4]{MurtySaradha2007}.
These results do not imply independence over the present $B$.

\subsection*{Detailed coefficient-cancellation argument}
Main-text Proposition~\ref{prop:coefficient-cancellation} is an algebraic
statement about the actual column $d$, independent of its special-function
origin. Write $d=d_0+M\xi$ with $d_0\in B^N$, $\rank_B M=\kappa$
and $1,\xi_1,\ldots,\xi_\kappa$ independent over $B$.
For $P\in\mathcal A_d$ of degree $m$, let $P_j$ be its homogeneous part
of degree $j$. The degree $m-1$ part of $P(Z-d)$ is
\[
 P_{m-1}(Z)-d_0\cdot\nabla P_m(Z)
 -\sum_j\xi_jM_j\cdot\nabla P_m(Z).
\]
Its coefficients lie in $B$, so each $M_j\cdot\nabla P_m=0$.
Thus $P_m$ is invariant under translation by $\operatorname{im}M$ and
$P_m(Z-d)=P_m(Z-d_0)\in B[Z]$. Subtract $P_m$ and induct on degree.
Conversely, this invariance makes $P(Z-d)=P(Z-d_0)\in B[Z]$.
Coordinates adapted to $\operatorname{im}M$ therefore give
$\mathcal A_d=B[JZ]$. Exact invariance additionally requires invariance
under $d_0$. Iteration at integer multiples, followed by polynomial
dependence on the translation parameter, gives invariance under $Bd_0$.
The joint invariant algebra is $B[J_0Z]$.

Since $Jd\in B^{N-\kappa}$, $B[JV]=B[JX]$ and its transcendence
degree is $N-\kappa$.
Since $d=A\mathbf c$ and $A$ is algebraic, $\kappa\le\rank_B A=\rank_{\Qbar}A$.
If the rows of a matrix $Q$ lie in
$\operatorname{row}(J)$, then $Qd\in B^s$ and
$\trdeg_B B(QV)=\rank_BQ$. Its relation ideal is generated by
the linear row dependencies $\mu W$, $\mu Q=0$: translation by $Qd$
preserves these equations because $\mu Qd=0$.
In particular $J_0V=J_0X$ and
$\trdeg_B B(J_0V)=\rank_BJ_0$.
The actual rank $\kappa$ and matrix $J$ are not claimed effective.
For $J_A$ whose rows span the left kernel $\{\lambda\in B^{1\times N}:\lambda A=0\}$, the uniform translation-invariant
polynomial and rational algebras are $B[J_AZ]$ and $B(J_AZ)$.
The actual row space $\{\lambda\in B^{1\times N}:\lambda d\in B\}$,
spanned by the rows of $J$, contains the left kernel of $A$ and may give
the larger coefficient-cancellation algebra $\mathcal A_d=B[JZ]$.
Once this row space is fixed, nonlinear relations among the entries of $d$
do not enlarge $\mathcal A_d$. Exact translation invariants are $B[J_0Z]$.

\subsection*{Computation for the mixed example}
In main-text Example~\ref{ex:ct2-mixed}, choose
$x=1$, $(u,m)=(2,0)$ and $(t,a)=(1,1/2)$.
Gauss's formula gives $G_{1/2}=-2\log2$; its cotangent term vanishes
and its sine sum is empty. All node logarithms and both recurrence
corrections vanish. Thus
\[
\begin{gathered}
 \mathbf c=\binom\gamma{\log2},\quad
 A=\begin{pmatrix}1&0\\1&0\\1&-2\end{pmatrix},\\
 X=\left(\Ein(1),\frac{\mathsf S(2)}{I_0(2)},
 \frac{S_{1/2}(1)}{f_{1/2}(1)}\right)^t,\quad V=X-A\mathbf c.
\end{gathered}
\]
No additive $\pi$ entry is needed, despite the multiplicative
$\Gamma(1/2)$ in the third connection value.
The first two constants cancel:
\[
 E_1(1)-\frac{K_0(2)}{I_0(2)}
 =\Ein(1)-\frac{\mathsf S(2)}{I_0(2)}=X_1-X_2.
\]
Main-text Theorem~\ref{thm:connection}, via Theorem~\ref{thm:values},
gives independence of $X_1,X_2$ over $B_*$. Thus this difference is
transcendental over $B_*$, hence over $\mathcal E$, since a polynomial
relation would contradict that independence.

\subsection*{What the actual-number bound asserts}
Main-text Theorem~\ref{thm:connection} gives
$\trdeg_F F(\mathbf c,V)\ge N$ for $F=\mathcal E,B$.
It assumes neither independence of $\mathbf c$ nor independence of $X$ over $F(\mathbf c)$.
For example, over an algebraically closed field $F$, take independent
$X_1,\ldots,X_N$, $\mathbf c=(X_1)$ and $d=(X_1,0,\ldots,0)$.
Then $V=(0,X_2,\ldots,X_N)$ and the combined degree is $N$,
although $\mathbf c$ is transcendental. Thus the argument does not give the
additive bound $N+\trdeg_F F(\mathbf c)$.
The mixed example is a simultaneous consequence of the joint value
theorem. Separate one-family statements alone do not yield independence
after all regular pairs have been adjoined. No priority claim is made.
